% !PRINTED 664
\weberpage{664}
%<<CONT
\gesperrt{so ordnen wir einem solchen Zusammentreffen von Werten
einen Punkt $\mathfrak{P}$ zu.} Wir sagen: $\alpha = \alpha_0$, oder $\alpha$ hat den Wert $\alpha_0$
im Punkte $\mathfrak{P}$. Zwei Punkte $\mathfrak{P}$ und $\mathfrak{P}'$ heißen dann und nur dann
verschieden, wenn es wenigstens eine Funktion in $\Omega$ gibt, die in
$\mathfrak{P}$ und $\mathfrak{P}'$ verschiedene Werte hat.

Die Regeln 2. bis 5. versagen in den Ausnahmefällen, nämlich
2. und 3., wenn $\alpha_0$ und $\beta_0 = \infty$ sind; 4., wenn $\alpha_0 = 0$, $\beta_0 = \infty$
oder $\alpha_0 = \infty$, $\beta_0 = 0$ wird, und 5., wenn $\alpha_0$ und $\beta_0$ beide $= 0$
oder beide $= \infty$ sind.

Trotzdem müssen die Funktionen $(\alpha \pm \beta)_0$, $(\alpha\beta)_0$ und $(\alpha : \beta)_0$
auch in diesen Fällen bestimmte Werte haben, die aber nicht unmittelbar
 aus den Vorschriften 2. bis 5., sondern auf indirektem
Wege bestimmt werden.

\gesperrt{Der Punkt ist hiernach ein zu dem Körper $\Omega$ gehöriger
 invarianter Begriff}, der nichts mit dem zufälligen
Umstand zu tun hat, welche der Funktionen von $\Omega$ wir als die
unabhängige Variable betrachten.

Um alle Punkte zu finden, verfahre man so: Ist $\mathfrak{P}$ ein Punkt,
so gibt es Funktionen, die in ihm einen endlichen Wert haben,
denn wenn eine Funktion $\alpha$ in $\mathfrak{P}$ unendlich ist, so ist $1 : \alpha$ endlich.
Eine solche Funktion nehme man als unabhängige Variable $z$
und bezeichne ihren Wert im Punkt $\mathfrak{P}$ mit $c$. Alle ganzen Funktionen
 von $z$ sind dann gleichfalls in $\mathfrak{P}$ endlich, denn ist $\omega$ eine
ganze Funktion, die der Gleichung
\[
\omega^{n} + a_1\omega^{n-1} + \cdots + a_n = 0
\]
genügt, so muß
\[
1 + \frac{a_1}{\omega} + \frac{a_2}{\omega^2} + \cdots + \frac{a_n}{\omega^n} = 0
\]
in $\mathfrak{P}$ befriedigt sein, und da die $a_1, a_2, \dots, a_n$ als ganze rationale
Funktionen von $z$ in $\mathfrak{P}$ endlich sind, so kann $\omega$ nicht $= \infty$ sein.

Jede Kongruenz nach dem Modul $z - c$ muß daher im
Punkt $\mathfrak{P}$ in eine richtige Gleichung übergehen, auch wenn diese
Kongruenz zwischen holomorphen ganzen Funktionalen besteht.
Ist also $\tau$ die Basisform von $\Omega$ in $Z$, so muß nach § 184, (20) die
Gleichung bestehen:
\[
(\tau - \tau_1)(\tau - \tau_2)\dots(\tau - \tau_n) = 0.
\]
Es muß also einer der Faktoren der linken Seite in $\mathfrak{P}$ verschwinden.
Ist dies $\tau - \tau_1$, so entspricht diesem Faktor ein
Primfaktor $\pi$ von $z - c$. Es wird also in $\mathfrak{P}$

\weq{2}{\omega_1 = c_1, \quad \omega_2 = c_2, \dots, \omega_n = c_n,}
%CONT>>
